% Einheit 1 von 3 – Summe mal Summe (bank/binomische-formeln/e1.jsonl, zusammenbau v0.1)
%% TODO Zweigzeile Teil 1: was hier gelernt wird (Satz in Schülersprache aus dem Einheitstitel) – Einheit 1
\einheitenkopf[e1]{Einheit 1 von 3 · Summe mal Summe}
\zweigzeile{ab Kl. 8, je nach Buch bis Kl. 9 (am Gymnasium ab Kl. 7) · keine P10-Aufgabe}
\setcounter{aufgabe}{9}

%% TODO Titel: Ich-kann-Satz fehlt in der Bank (Platzhalter: Kettenname) – Nr. 10
%% TODO Anweisung: ein Satz über der ersten Teilaufgabe fehlt in der Bank – Nr. 10
\begin{aufgabe}{Summe mal Summe}
%% TODO \rechenplatz{n}: Schrittzahl des Grundfalls fehlt in der Bank (nur bei fester Schreibform, 2.3 b)
\begin{teile}
\teil Verbinde jedes Glied der ersten Klammer mit jedem Glied der zweiten durch einen Pfeil und schreib die Produkte hin, ohne zu rechnen: $(x + 4) \cdot (y + 6)$ \leerfeld
\teil Fasse zusammen: $3x + 4y + 2x + y$ \leerfeld
\teil Fasse zusammen: $6x - 2y - x + 5y$ \leerfeld
\teil Fasse zusammen: $2x^2 + 3xy + x^2 + 4x - xy$ \leerfeld
\end{teile}
\begin{gleichungsraster}[2]
\gl{\text{Multipliziere aus: $4 \cdot (2x + 3y)$}} &
\gl{\text{Schreib die vier Produkte hin: $(x + 4) \cdot (x + 7)$}} \\
\gl{\text{Multipliziere aus: $(x + 2) \cdot (x + 8)$}} &
\gl{\text{Multipliziere aus: $(x + 6) \cdot (x - 2)$}} \\
\gl{\text{Multipliziere aus: $(x - 6) \cdot (x - 1)$}} &
\gl{\text{Multipliziere aus: $(2x + 1) \cdot (x + 4)$}} \\
\gl{\text{Multipliziere aus: $(a + 2b) \cdot (3a + b)$}} &
\gl{\text{Multipliziere aus und fasse zusammen: $(3x - 4) \cdot (2x + 5)$}} \\
\end{gleichungsraster}
\end{aufgabe}

%% TODO Titel: Ich-kann-Satz fehlt in der Bank (Platzhalter: Kettenname) – Nr. 11
%% TODO Anweisung: ein Satz über der ersten Teilaufgabe fehlt in der Bank – Nr. 11
\begin{aufgabe}{Termwert mit zwei Variablen berechnen (auch negativ)}
\begin{teile}
\teil $3x + 2y$ für $x = 4$ und $y = -5$ – Termwert? \leerfeld
\end{teile}
\end{aufgabe}

%% TODO Titel: Ich-kann-Satz fehlt in der Bank (Platzhalter: Kettenname) – Nr. 12
%% TODO Anweisung: ein Satz über der ersten Teilaufgabe fehlt in der Bank – Nr. 12
\begin{aufgabe}{Summe mal Summe – Fehler finden, Begründen}
\begin{teile}
\teil Finde den Fehler und rechne richtig: \rechnung{(x + 6) \cdot (x + 4) &= x^2 + 24}
\teil Warum entstehen bei $(x + 2) \cdot (y + 5)$ vier Produkte? Begründe am Flächenbild eines Rechtecks.
\end{teile}
\end{aufgabe}
